\documentclass[11pt]{article}
\usepackage{icmlko}
\begin{document}

\title{드리프트는 있는데 살 수가 없다 --- 예보 비용을 창으로도 무게로도
못 줄인다}
\author{월드모델 실험실}
\date{노트 451 · 2026년 8월 2일}
\maketitle

\begin{abstract}
노트 445 가 이 실험실에서 제일 큰 수를 찾아 놓고 손을 못 대고 있었다 ---
\textbf{예보 값}(혼자 예보 $-$ 같은 시기 천장)이 평균 $-0.1492$ 로, 모음이
사는 $+0.0699$ 보다 크고 열 도메인 중 일곱에서 더 크다. 노트 447 의
정직한 시간 축이 산 것은 $+0.004$ 뿐이었다. 그래서 먼저 물었다 ---
\textbf{이 비용은 줄어드는 종류인가.}

가장 흔한 답부터 쟀다. \emph{옛 행이 상해서}라면 최근 창으로만 배우면
낫다. \textbf{전부 내려간다.}

\begin{center}
\begin{tabular}{lrr}
\hline
학습 창 & 학습 행 & 판 차 \\
\hline
전체 & $18{,}540$ & --- \\
2018$+$ & $11{,}320$ & $-0.0017$ \\
2020$+$ & $8{,}455$ & $-0.0099$ \\
2022$+$ & $5{,}398$ & $-0.0308$ \\
\hline
\end{tabular}
\end{center}

\textbf{여기서 멈추면 틀린다}(노트 133). 자르기는 \emph{상한 행을 빼는
것}과 \emph{행을 빼는 것}을 한꺼번에 한다. 드리프트가 진짜라도 행 수
손해가 더 크면 자르기는 진다.

그래서 \textbf{행 수를 똑같이 고정}했다. $N = 8{,}455$ 로 맞추고
\emph{최근 $N$} 과 \emph{전 기간 무작위 $N$} 을 견줬다.

\begin{center}
\begin{tabular}{lrr}
\hline
 & 판 & \\
\hline
전체 $18{,}540$행 & $+0.4860$ & \\
무작위 $8{,}455$행 & $+0.4699$ & 크기 $\mathbf{-0.0161}$ \\
\textbf{최근 $8{,}455$행} & $+0.4755$ & 최근성 $\mathbf{+0.0057}$ \\
\hline
\end{tabular}
\end{center}

\textbf{드리프트는 실재한다} --- $+0.0057$, SE $0.0021$, 95\%
$[+0.0016, +0.0097]$, 구간이 $0$ 을 안 문다. 날짜 통제 차도 $+0.0054$ 라
노트 446 의 구멍도 아니다. 그리고 자르기의 손해가 정확히 갈린다:
$\mathbf{-0.0099 = -0.0161 + 0.0057}$.

\textbf{그러면 자르지 말고 무게를 주면 된다} --- 고 생각했다.
\textbf{반증됐다.}
\end{abstract}

\begin{figure}[t]
\centering
\includegraphics[width=\textwidth]{figs/f1.pdf}
\caption{\textbf{가} --- 행을 버리는 길(자르기)과 행을 덜 보는
길(무게)이 나란히 내려간다. \textbf{나} --- 행 수를 고정하면 최근성이
보인다. 막대의 오차는 뽑기 열 번의 95\% 구간.}
\end{figure}

\section{무게로는 못 산다}

전 행을 그대로 쓰되 나이로 지수 감쇠 무게를 줬다(반감 $\tau$).

\begin{center}
\begin{tabular}{lrr}
\hline
$\tau$ & 판 차 & 날짜 통제 차 \\
\hline
$8$년 & $-0.0005$ & $-0.0013$ \\
$4$년 & $-0.0038$ & $-0.0039$ \\
$2$년 & $-0.0130$ & $-0.0141$ \\
$1$년 & $-0.0294$ & $-0.0328$ \\
\hline
\end{tabular}
\end{center}

\textbf{어떤 $\tau$ 도 균등을 못 이긴다.} $\tau$ 를 줄일수록 단조로
나빠지고, 제일 무해한 $\tau=8$ 조차 $-0.0005$ 다.

등록할 때 적어 둔 반증 방식이 그대로 나왔다 --- \emph{최근성의 값은 행의
\textbf{무게}가 아니라 행의 \textbf{조합}에 있다.} $8{,}455$행만 가질 수
있을 때는 최근 것으로 채우는 게 낫지만, $18{,}540$행을 다 가진 뒤에는
옛 행을 덜 보는 것이 아무것도 안 사고 유효 표본만 깎는다.

\section{자기 정정 --- 첫 답은 SE 한 개였다}

씨앗 넷으로 처음 쟀을 때 최근 $-$ 무작위가 $+0.0032$ 였고, 미리 적어 둔
문턱 $+0.003$ 을 아슬하게 넘어 ``통과''가 찍혔다. \textbf{그 통과는
잡음이었다.} 무작위 뽑기 넷이 $+0.4689 \cdot +0.4722 \cdot +0.4796 \cdot
+0.4643$ 으로 SD $0.0063$ 이니 넷 평균의 SE 가 $0.0032$ --- \textbf{차이가
정확히 SE 한 개}였다.

문턱을 잡을 때 \emph{흩어짐을 모르고} 잡은 것이 잘못이다. 뽑기를 열로
늘려 구간으로 다시 정하니 $[+0.0016, +0.0097]$ 로 살아남았다. 결론은
같지만 \textbf{처음 근거로는 그 결론을 말할 수 없었다} --- 노트 434 가
LODO 에서 겪은 것과 같은 병이다.

\section*{규약 --- 문턱은 흩어짐을 재고 나서 잡는다}

사전 등록은 좋은 습관이지만 \textbf{문턱을 아무 데나 잡으면 등록이
잡음을 통과시킨다}. 이 노트가 붙이는 조항: \textbf{뽑기가 들어가는
비교는 문턱을 잡기 전에 뽑기 흩어짐을 먼저 재고, 문턱을 그 SE 의 두 배
위에 잡는다.} 못 재고 시작했으면 결과를 ``통과''로 적지 말고 흩어짐을
재서 다시 정한다.

\section*{남은 것}

\begin{itemize}
\item \textbf{$-0.1492$ 는 그대로 남았다.} 창도 무게도 아니라면 남은
후보는 \emph{예보 자체를 다르게 하는 것}이다 --- 유보 시점에 이미 보이는
초기 신호(첫 주 추세 · 사전 예약)를 재료로 쓰는 것. 지금 축은 전부
출시 전 정보다.
\item \textbf{드리프트가 어디서 오나.} $+0.0057$ 이 도메인마다 다를 텐데
안 갈라 봤다. 웹툰처럼 형식이 빨리 변하는 데가 크면 그 도메인만
최근 창을 줄 수도 있다.
\item \textbf{빈 칸 네 도메인}(노트 450) --- 문턱으로는 못 채운다.
\end{itemize}

\end{document}
